% Source: https://artofproblemsolving.com/community/c3h2601080_trivial_fact
% Stable identity: aops-topic-2601080-q01
Let the circle have centre $O$ and radius $r$, and let $AB$ be any chord. By the triangle inequality (or trivially, if $O$ lies on line $AB$),
\[AB\le AO+OB=r+r=2r,\]
with equality exactly when $O$ lies on the segment $AB$, that is, when $AB$ is a diameter. So every chord has length at most $2r$, and the chords of length $2r$ are exactly the diameters. $\blacksquare$

\begin{asy}
usepackage("amsmath");
size(3.5cm);
defaultpen(fontsize(10pt));
pair O=(0,0),A=dir(200),B=dir(-35),P=dir(100),Q=dir(-80);
draw(unitcircle,linewidth(0.8));draw(A--B,linewidth(0.8));draw(A--O--B,dashed);draw(P--Q,dotted);
dot(O);dot(A);dot(B);
label("$O$",O,NE);label("$A$",A,W);label("$B$",B,SE);label("$r$",(A+O)/2,N,fontsize(8pt));label("$r$",(B+O)/2,NE,fontsize(8pt));
\end{asy}

The dotted segment is a diameter, of length $2r$.
